\documentclass[tikz,border=6pt]{standalone}
\usepackage[no-math]{fontspec}
\usepackage[brazil]{babel}
\usepackage{amsmath,amssymb}
\IfFileExists{C:/Windows/Fonts/Lato_Regular.ttf}{%
  \setmainfont{Lato}[
    Path = C:/Windows/Fonts/,
    UprightFont = Lato_Regular.ttf,
    BoldFont = Lato_Bold.ttf,
    ItalicFont = Lato_Italic.ttf,
    BoldItalicFont = Lato_BoldItalic.ttf
  ]
  \setsansfont{Lato}[
    Path = C:/Windows/Fonts/,
    UprightFont = Lato_Regular.ttf,
    BoldFont = Lato_Bold.ttf,
    ItalicFont = Lato_Italic.ttf,
    BoldItalicFont = Lato_BoldItalic.ttf
  ]
}{%
  \RequirePackage[default]{lato}
}
\renewcommand{\familydefault}{\sfdefault}

\usepackage{xcolor}
\usepackage{tikz}
\usepackage{pgfplots}
\pgfplotsset{compat=1.18}
\usetikzlibrary{arrows.meta,positioning,shapes.geometric,fit,backgrounds,calc,decorations.pathreplacing,patterns}

% Paleta institucional UEFS / CHF662
\definecolor{uefsNavy}{HTML}{1B2A4A}
\definecolor{uefsTeal}{HTML}{0E7C7B}
\definecolor{uefsForest}{HTML}{2D6A4F}
\definecolor{uefsGold}{HTML}{C4983B}
\definecolor{uefsGoldDark}{HTML}{8F6B1E}
\definecolor{uefsAlert}{HTML}{B02A37}
\definecolor{uefsOrange}{HTML}{D95F02}
\definecolor{uefsInk}{HTML}{22252A}
\definecolor{uefsMuted}{HTML}{5C6470}
\definecolor{uefsWater}{HTML}{2E86AB}
\definecolor{uefsSand}{HTML}{D9B382}
\definecolor{bgNavy}{HTML}{F2F5F9}
\definecolor{bgTeal}{HTML}{E9F5F4}
\definecolor{bgForest}{HTML}{EEF5F0}
\definecolor{bgGold}{HTML}{FCF7EB}
\definecolor{bgAlert}{HTML}{FDF2F2}
\definecolor{bgGray}{HTML}{F8FAFC}
\definecolor{lineGray}{HTML}{CFD8DC}

\tikzset{
  >=Stealth,
  aaCard/.style={draw=#1, fill=white, rounded corners=4pt, line width=1pt, inner sep=6pt, align=center, font=\small},
  aaBox/.style={draw=#1, rounded corners=4pt, line width=1.1pt, inner sep=6pt, align=center, font=\footnotesize\sffamily, text=uefsInk},
  aaTitle/.style={font=\normalsize\bfseries\sffamily, text=uefsNavy, align=center},
  aaSub/.style={font=\footnotesize\sffamily, text=uefsMuted, align=center},
  aaArrow/.style={->, >=Stealth, line width=1.4pt},
  aaDArrow/.style={<->, >=Stealth, line width=1.4pt}
}


\begin{document}
\begin{tikzpicture}
\begin{axis}[
  width=15.5cm,
  height=8.2cm,
  xmin=0, xmax=10,
  ymin=0, ymax=5.2,
  xtick={0,2,4,6,8,10},
  ytick={0,1,2,3,4,5},
  xlabel={Perfil longitudinal do canal erosivo (montante $\rightarrow$ jusante)},
  ylabel={Cota relativa (m)},
  xlabel style={font=\small\sffamily, text=uefsInk},
  ylabel style={font=\small\sffamily, text=uefsInk},
  title={Bioengenharia de Água e Solo: Estabilização de Ravina por Paliçadas Permeáveis em Série},
  title style={font=\normalsize\bfseries\sffamily, text=uefsNavy, yshift=4pt},
  axis lines*=left,
  axis line style={color=lineGray, line width=0.9pt},
  tick style={color=uefsMuted},
  xticklabel style={font=\footnotesize\sffamily, text=uefsInk},
  yticklabel style={font=\footnotesize\sffamily, text=uefsInk},
  clip=false
]

% Solo original abaixo do perfil íngreme (y = 3,6 - (x - 0,5)/3)
\fill[uefsSand!45!white]
  (axis cs:0.5, 0) -- (axis cs:0.5, 3.60) -- (axis cs:9.5, 0.60) -- (axis cs:9.5, 0) -- cycle;

% 3 cunhas escalonadas de sedimento retido
% Cunha 1: x=1,0 (y=3,433) até x=2,92 (y=3,38)
\fill[uefsSand!90!uefsGoldDark]
  (axis cs:1.0, 3.433) -- (axis cs:2.92, 3.38) -- (axis cs:2.92, 2.793) -- cycle;
\draw[uefsNavy, line width=0.8pt] (axis cs:1.0, 3.433) -- (axis cs:2.92, 3.38);

% Cunha 2: x=3,5 (y=2,60) até x=5,42 (y=2,54)
\fill[uefsSand!90!uefsGoldDark]
  (axis cs:3.5, 2.60) -- (axis cs:5.42, 2.54) -- (axis cs:5.42, 1.96) -- cycle;
\draw[uefsNavy, line width=0.8pt] (axis cs:3.5, 2.60) -- (axis cs:5.42, 2.54);

% Cunha 3: x=6,0 (y=1,767) até x=7,92 (y=1,71)
\fill[uefsSand!90!uefsGoldDark]
  (axis cs:6.0, 1.767) -- (axis cs:7.92, 1.71) -- (axis cs:7.92, 1.127) -- cycle;
\draw[uefsNavy, line width=0.8pt] (axis cs:6.0, 1.767) -- (axis cs:7.92, 1.71);

% Linha tracejada do perfil original íngreme
\draw[dashed, color=uefsAlert, line width=1.6pt]
  (axis cs:0.5, 3.60) -- (axis cs:9.5, 0.60);

% 3 paliçadas permeáveis em série (madeira marrom com borda navy)
\draw[draw=uefsNavy, fill=uefsGoldDark!55!black, line width=1pt]
  (axis cs:2.92, 2.38) rectangle (axis cs:3.10, 3.46);
\draw[draw=uefsNavy, fill=uefsGoldDark!55!black, line width=1pt]
  (axis cs:5.42, 1.56) rectangle (axis cs:5.60, 2.62);
\draw[draw=uefsNavy, fill=uefsGoldDark!55!black, line width=1pt]
  (axis cs:7.92, 0.74) rectangle (axis cs:8.10, 1.79);

% Estacas/mudas verdes posicionadas BEM ABAIXO das caixas de chamada superiores
% Cunha 1 (superfície ~3,40 -> topo das mudas ~3,82; caixa superior fica acima de y=4,35)
\draw[color=uefsForest, line width=1.6pt] (axis cs:1.75, 3.11) -- (axis cs:1.75, 3.77);
\fill[uefsForest] (axis cs:1.75, 3.83) circle (4.5pt);
\draw[color=uefsForest, line width=1.6pt] (axis cs:2.35, 3.10) -- (axis cs:2.35, 3.76);
\fill[uefsForest] (axis cs:2.35, 3.82) circle (4.5pt);
% Cunha 2 (superfície ~2,56 -> topo das mudas ~2,98)
\draw[color=uefsForest, line width=1.6pt] (axis cs:4.35, 2.29) -- (axis cs:4.35, 2.93);
\fill[uefsForest] (axis cs:4.35, 2.99) circle (4.5pt);
\draw[color=uefsForest, line width=1.6pt] (axis cs:4.90, 2.27) -- (axis cs:4.90, 2.91);
\fill[uefsForest] (axis cs:4.90, 2.97) circle (4.5pt);
% Cunha 3 (superfície ~1,73 -> topo das mudas ~2,15)
\draw[color=uefsForest, line width=1.6pt] (axis cs:6.80, 1.46) -- (axis cs:6.80, 2.10);
\fill[uefsForest] (axis cs:6.80, 2.16) circle (4.5pt);
\draw[color=uefsForest, line width=1.6pt] (axis cs:7.35, 1.44) -- (axis cs:7.35, 2.08);
\fill[uefsForest] (axis cs:7.35, 2.14) circle (4.5pt);

% Caixas explicativas no espaço livre superior (sem qualquer contato com mudas)
\node[
  draw=uefsForest, fill=bgForest, rounded corners=4pt, line width=1.1pt,
  inner sep=5pt, align=left, font=\footnotesize\bfseries\sffamily, text=uefsForest
] (boxCunha) at (axis cs:2.25, 4.65)
  {Cunha de sedimento + revegetação\\(novo perfil escalonado mais suave)};

\node[
  draw=uefsNavy, fill=bgNavy, rounded corners=4pt, line width=1.1pt,
  inner sep=5pt, align=left, font=\footnotesize\bfseries\sffamily, text=uefsNavy
] (boxPalicada) at (axis cs:7.25, 4.35)
  {Paliçada permeável\\(reduz velocidade do fluxo,\\retém sedimentos, drena água)};

% Setas limpas das caixas para a cunha 2 (trecho livre x=3,85) e paliçada 2 (x=5,55)
\draw[aaArrow, color=uefsForest, line width=1.4pt]
  (boxCunha.south east) -- (axis cs:3.85, 2.64);
\draw[aaArrow, color=uefsNavy, line width=1.4pt]
  (boxPalicada.south west) -- (axis cs:5.62, 2.65);

% Escoamento dissipado à jusante
\draw[aaArrow, color=uefsWater, line width=1.6pt]
  (axis cs:8.28, 1.32) .. controls (axis cs:8.72, 1.16) and (axis cs:9.05, 1.05) .. (axis cs:9.45, 0.90);

\node[
  draw=uefsWater, fill=white, rounded corners=3pt, line width=0.9pt,
  inner sep=4pt, align=center, font=\footnotesize\bfseries\sffamily, text=uefsWater
] (boxEsc) at (axis cs:8.65, 2.35)
  {Escoamento dissipado\\(menor energia cinética)};
\draw[aaArrow, color=uefsWater, line width=1.2pt]
  (boxEsc.south) -- (axis cs:8.85, 1.22);

% Legenda inferior esquerda
\node[
  draw=lineGray, fill=white, rounded corners=3pt, line width=0.8pt,
  inner sep=4pt, anchor=south west, font=\footnotesize\sffamily, text=uefsInk
] at (axis cs:0.22, 0.22)
  {\tikz[baseline=-0.5ex]{\draw[dashed, color=uefsAlert, line width=1.5pt] (0,0) -- (0.65,0);}~Perfil original íngreme do leito da ravina};

\end{axis}
\end{tikzpicture}
\end{document}
